\chapter{Methodik}
%
\section{Forschungsfrage und experimenteller Rahmen}
%
Die zentrale Forschungsfrage dieser Arbeit lautet:
Welche semantischen Eigenschaften eines Float-Embedding-Raums bleiben unter Quantisierung erhalten, welche gehen verloren, und wie wirkt sich dieser Verlust auf die Retrieval-Qualität aus?
%
Diese Formulierung ist bewusst enger gefasst als eine allgemeine Frage nach der Nützlichkeit von Bit-Embeddings.
Sie verlangt nicht nur die Messung eines Qualitätsverlusts, sondern auch dessen Erklärung aus den Eigenschaften des Ausgangsraums.
Damit geht die Arbeit über einen reinen Benchmark-Vergleich hinaus und liefert ein Verständnis dafür, \emph{warum} bestimmte Kompressionsverfahren auf bestimmten Daten besser oder schlechter funktionieren.
%
Um diese Frage empirisch zu beantworten, wird ein kontrolliertes Experiment durchgeführt, das sich in fünf Phasen gliedert:
\begin{enumerate}
  \item \textbf{Charakterisierung:} Der Float-Raum wird statistisch und geometrisch analysiert, um Erwartungen an den Quantisierungsverlust zu formulieren.
  \item \textbf{Kompression:} Aus dem Float-Raum werden komprimierte Repräsentationen abgeleitet -- variiert entlang zweier Achsen (Bittiefe und Dimensionszahl).
  \item \textbf{Strukturanalyse:} Die Fähigkeit der komprimierten Räume zur Distanz- und Nachbarschaftserhaltung wird gemessen.
  \item \textbf{Retrieval-Evaluation:} Die tatsächliche Auswirkung auf die Suchqualität wird mit etablierten Information-Retrieval-Metriken quantifiziert.
  \item \textbf{Praktische Evaluation:} Speicherbedarf und Suchgeschwindigkeit werden gemessen, um Qualität und Effizienz gegeneinander abzuwägen.
\end{enumerate}
%
Diese Reihenfolge ist methodisch notwendig.
Würde direkt mit der Retrieval-Evaluation begonnen, ließe sich zwar feststellen, \emph{dass} ein Verfahren schlechter abschneidet, aber nicht \emph{warum}.
Die vorgeschaltete Raumanalyse ermöglicht es, Qualitätsverluste kausal auf geometrische Eigenschaften des Ausgangsraums zurückzuführen.
%
\section{Datensatz}
%
Als Evaluationsgrundlage dient der BEIR-Benchmark \cite{thakur_2021_beir}.
BEIR stellt annotierte Relevanzurteile zwischen Anfragen und Dokumenten bereit und deckt heterogene Domänen ab.
Die Wahl fällt auf BEIR aus zwei Gründen:
Erstens liegen manuell geprüfte Relevanzannotationen vor, sodass Retrieval-Qualität objektiv messbar ist und nicht auf synthetischen Ähnlichkeitsmaßen beruht.
Zweitens erlaubt die Domänenvielfalt eine Aussage über die Generalisierbarkeit der Ergebnisse.
%
Innerhalb von BEIR werden drei Subsets verwendet:
\begin{itemize}
  \item \textbf{SciFact} -- wissenschaftliche Faktenverifikation mit kurzen Dokumenten und binären Relevanzurteilen.
  \item \textbf{FiQA} -- Frage-Antwort-Paare aus dem Finanzbereich mit längeren, argumentativen Texten.
  \item \textbf{TREC-COVID} -- biomedizinische Suchanfragen mit umfangreichen Dokumenten und mehrstufigen Relevanzannotationen.
\end{itemize}
%
Die Auswahl folgt dem Prinzip maximaler Variation entlang dreier Achsen:
Dokumentlänge (kurz, mittel, lang), Fachsprache (Naturwissenschaft, Finanz, Biomedizin) und Annotationsschema (binär vs. mehrstufig).
Damit kann systematisch geprüft werden, ob Quantisierungseffekte domänenspezifisch auftreten oder sich über verschiedene Texttypen hinweg konsistent verhalten.
Würde nur ein einzelnes Subset verwendet, wäre unklar, ob die Ergebnisse an der Methode oder an den Daten liegen.
%
\section{Embedding-Modell}
%
Für die Erzeugung aller Vektorrepräsentationen wird ausschließlich das Modell \texttt{bge-large-en-v1.5} \cite{xiao_2023_bge} verwendet.
Die Beschränkung auf ein einziges Modell ist methodisch notwendig, um Konfundierung zu vermeiden:
Unterschiedliche Modelle erzeugen Vektorräume mit unterschiedlicher Geometrie, sodass ein Vergleich der Quantisierungsergebnisse nicht mehr eindeutig auf die Kompressionsverfahren zurückführbar wäre.
%
Die Wahl fällt auf \texttt{bge-large-en-v1.5} aus folgenden Gründen:
\begin{itemize}
  \item Das Modell erreicht zum Zeitpunkt der Modellauswahl konsistent hohe Werte auf dem MTEB-Leaderboard \cite{muennighoff_2022_mteb} und gilt als etablierter Referenzpunkt in der Embedding-Literatur.
  \item Es erzeugt Vektoren mit 768 Dimensionen -- genügend Struktur für eine aussagekräftige Quantisierungsanalyse, aber nicht so hochdimensional, dass die Berechnungen unpraktikabel werden.
  \item Als Open-Source-Modell ermöglicht es vollständige Reproduzierbarkeit.
\end{itemize}
%
Eine naheliegende Alternative wäre \texttt{nomic-embed-text-v1.5} \cite{nussbaum_2024_nomic}, das durch Matryoshka-Training bereits auf variable Dimensionsgrößen und binäre Ausgabe optimiert ist.
Genau diese Optimierung macht es für die vorliegende Untersuchung jedoch weniger geeignet:
Da das Ziel ist, den Informationsverlust bei nachträglicher Quantisierung zu messen, würde ein bereits auf Kompression trainiertes Modell den Effekt der Quantisierungsverfahren selbst verschleiern.
Die Ergebnisse wären nicht mehr auf das Kompressionsverfahren, sondern auf die Trainingsoptimierung des Modells zurückführbar.
%
\section{Erzeugung der Repräsentationsräume}
%
Ausgangspunkt ist der vollständige Float-Raum.
Für jedes Dokument im Datensatz wird ein Embedding-Vektor $\mathbf{x} \in \mathbb{R}^{768}$ berechnet.
Die resultierende Matrix $\mathbf{X} \in \mathbb{R}^{N \times 768}$ bildet die Ground Truth, gegen die alle komprimierten Repräsentationen evaluiert werden.
%
Aus diesem Float-Raum werden zwei komprimierte Varianten \emph{nachträglich} abgeleitet.
Dieser Ansatz -- nachträgliche Kompression eines bestehenden Float-Raums -- ist eine bewusste methodische Entscheidung gegenüber der Alternative, ein Modell zu verwenden, das direkt binäre oder quantisierte Ausgaben erzeugt (etwa durch Matryoshka-Training oder spezialisierte Binary-Embedding-Modelle).
Die nachträgliche Kompression ist für die Forschungsfrage aus drei Gründen notwendig:
\begin{enumerate}
  \item Sie ermöglicht einen kontrollierten Vorher-Nachher-Vergleich. Ohne einen Float-Referenzraum kann der Informationsverlust nicht gemessen werden, weil kein unkomprimierter Zustand existiert, gegen den evaluiert werden könnte.
  \item Sie isoliert genau eine unabhängige Variable -- das Kompressionsverfahren. Bei end-to-end trainierten Modellen wären Unterschiede in der Retrieval-Qualität nicht eindeutig auf die Kompression zurückführbar, sondern könnten ebenso an Architektur, Trainingsdaten oder Verlustfunktion liegen.
  \item Sie bildet das praxisrelevante Szenario ab, in dem bestehende Float-Embeddings nachträglich komprimiert werden sollen, ohne das zugrunde liegende Modell austauschen zu müssen.
\end{enumerate}
%
Die Wahl genau zweier komprimierter Varianten ist nicht willkürlich, sondern bildet die beiden Endpunkte eines Kompression-Qualität-Spektrums ab:
maximale Kompression bei minimaler Methodik (naive Binarisierung) und moderate Kompression bei informierter Methodik (TurboQuant).
%
\subsection{Binäre Repräsentation}
%
Die naive Binarisierung wendet einen Vorzeichenschwellwert an:
\begin{equation}
  b_i = \begin{cases} 1 & \text{falls } x_i > 0 \\ 0 & \text{sonst} \end{cases}
\end{equation}
%
Jede Dimension wird auf ein einzelnes Bit reduziert, was einem Kompressionsfaktor von 32 gegenüber Float32 entspricht.
Die Ähnlichkeit zweier binärer Vektoren wird über die Hamming-Distanz bestimmt.
%
Dieses Verfahren dient als untere Basislinie:
Es ist das einfachste denkbare Quantisierungsverfahren, das keine Kenntnis über die Datenverteilung nutzt.
Jede informiertere Methode sollte besser abschneiden.
Falls die naive Binarisierung bereits gute Ergebnisse liefert, deutet dies darauf hin, dass der Float-Raum Eigenschaften besitzt, die Quantisierung inhärent begünstigen -- etwa eine Konzentration der Information im Vorzeichen.
%
\subsection{TurboQuant-Repräsentation}
%
Die zweite Variante nutzt das TurboQuant-Verfahren \cite{zandieh_2025_turboquant}, das eine datenunabhängige Rotation mit anschließender skalarer Quantisierung auf 2~Bit pro Dimension kombiniert.
Der Kompressionsfaktor beträgt 16 gegenüber Float32.
%
TurboQuant ist aus mehreren Gründen als Vergleichsmethode geeignet:
\begin{itemize}
  \item Es repräsentiert den aktuellen Stand der Forschung für datenunabhängige Vektorquantisierung.
  \item Im Gegensatz zur naiven Binarisierung nutzt es eine vorgeschaltete Rotation, die stark asymmetrische Dimensionen ausgleicht (vgl. Kapitel~\ref{chap:grundlagen}).
  \item Der Kompressionsfaktor liegt zwischen Float und Binär, was eine Einordnung entlang des Kompression-Qualität-Spektrums ermöglicht.
\end{itemize}
%
Der methodische Mehrwert des Vergleichs liegt darin, dass sich zwei Effekte trennen lassen:
der Effekt der Bitreduktion selbst (sichtbar am Unterschied zwischen Float und TurboQuant) und der Effekt fehlender Vorverarbeitung (sichtbar am Unterschied zwischen TurboQuant und naiver Binarisierung).
%
\subsection{Dimensionsreduktion als zusätzlicher Kompressionsparameter}
\label{sec:methodik_dimreduktion}
%
Die bisherigen Repräsentationen variieren ausschließlich die Bittiefe pro Dimension bei fester Dimensionszahl ($d = 768$).
Dies bildet jedoch nur eine Achse des Kompressionsraums ab.
Die zweite Achse -- die Anzahl der verwendeten Dimensionen -- wird ergänzend untersucht, um den vollständigen Trade-off zwischen Kompression und Qualität zu erfassen.
%
Konkret werden alle Quantisierungsverfahren zusätzlich auf dimensionsreduzierte Vektoren angewandt.
Die Reduktion erfolgt über eine PCA-Projektion auf die Hauptkomponenten mit der höchsten erklärten Varianz.
Evaluiert werden die Präfixlängen $d \in \{64, 128, 256, 384, 768\}$.
%
Daraus ergibt sich eine zweidimensionale Versuchsmatrix:
\begin{itemize}
  \item Bittiefe: Float32, 2~Bit (TurboQuant), 1~Bit (binär)
  \item Dimensionszahl: 64, 128, 256, 384, 768
\end{itemize}
%
Für jede Kombination werden dieselben Metriken erhoben wie für die Vollversionen (Distanzkorrelation, Neighborhood Overlap, NDCG@10).
Die Ergebnisse werden als Pareto-Front dargestellt, die den optimalen Trade-off zwischen Gesamtkompression (Bits pro Vektor) und Retrieval-Qualität zeigt.
%
Der Gesamtspeicherbedarf einer Kombination beträgt $d \times b$ Bit pro Vektor, wobei $b$ die Bittiefe ist.
Damit reicht das Spektrum von $768 \times 32 = 24\,576$~Bit (Float, volle Dimension) bis $64 \times 1 = 64$~Bit (binär, minimale Dimension) -- ein Kompressionsfaktor von bis zu 384.
%
Die Wahl von PCA als Reduktionsverfahren ist methodisch konsistent mit dem Gesamtansatz der Arbeit:
Sie ist datenunabhängig im Sinne, dass keine aufgabenspezifische Optimierung stattfindet, und sie ist orthogonal zur anschließenden Quantisierung, sodass beide Effekte weiterhin separierbar bleiben.
Zudem liefert die in Phase~1 ohnehin durchgeführte PCA-Analyse der erklärten Varianz direkt eine Erwartung darüber, ab welcher Präfixlänge substantielle Information verloren geht.
%
\section{Phase~1: Charakterisierung des Float-Raums}
\label{sec:methodik_phase1}
%
Bevor Quantisierung stattfindet, wird der Float-Raum analysiert.
Diese Phase dient nicht der Bewertung der Kompressionsverfahren, sondern der Formulierung von Erwartungen:
Aus den Eigenschaften des Raums lässt sich vorhersagen, wo Quantisierung gut funktionieren wird und wo nicht.
Diese Vorhersagen werden in den späteren Phasen empirisch überprüft.
%
\subsection{Normverteilung}
%
Für alle Vektoren wird die L2-Norm $\|\mathbf{x}\|_2$ berechnet und als Histogramm dargestellt.
Die Normverteilung beantwortet eine konkrete Frage:
Liegen die Vektoren annähernd auf einer Hypersphäre, oder existieren signifikante Normunterschiede?
%
Diese Analyse ist für die Binarisierung direkt relevant:
Die Vorzeichenfunktion verwirft die Norminformation vollständig.
Wenn die Normen stark variieren, geht bei der Binarisierung Information verloren, die nicht im Vorzeichen kodiert ist.
Wenn die Normen hingegen nahezu konstant sind, ist der Informationsgehalt der Norm gering und ihr Verlust unproblematisch.
%
\subsection{Dimensionsstatistiken}
%
Für jede der 768 Dimensionen werden vier Kennzahlen berechnet:
\begin{itemize}
  \item \textbf{Mittelwert und Standardabweichung} -- beschreiben Lage und Streuung.
  \item \textbf{Schiefe (Skewness)} -- zeigt, ob die Verteilung symmetrisch um Null liegt. Bei starker Schiefe liegt der Vorzeichenschwellwert $x_i = 0$ der Binarisierung nicht am informationstheoretisch optimalen Punkt, da eine Seite der Verteilung mehr Masse trägt und das resultierende Bit damit weniger Information kodiert.
  \item \textbf{Kurtosis} -- misst die Schwere der Verteilungsränder. Hohe Kurtosis deutet auf Ausreißer hin, die bei Quantisierung auf wenige Stufen besonders stark verzerrt werden, weil die extremen Werte auf dieselben Quantisierungsgrenzen abgebildet werden wie die häufigen.
\end{itemize}
%
Diese Kennzahlen sind nicht Selbstzweck, sondern liefern spezifische Vorhersagen:
Dimensionen mit hoher Schiefe sollten bei der Binarisierung mehr Fehler verursachen als symmetrische Dimensionen.
Dimensionen mit hoher Kurtosis sollten bei TurboQuant (2~Bit, also 4 Stufen) mehr Quantisierungsfehler zeigen.
Ob diese Vorhersagen zutreffen, wird in Phase~2 überprüft.
%
\subsection{Intrinsische Dimensionalität}
%
Die intrinsische Dimensionalität (vgl. Kapitel~\ref{chap:grundlagen}) wird über zwei komplementäre Verfahren geschätzt:
\begin{itemize}
  \item \textbf{TwoNN} \cite{facco_2017_twonn} -- ein nächste-Nachbarn-basierter Schätzer, der die Verhältnisse der Abstände zum ersten und zweiten Nachbarn auswertet. Er erfasst die lokale Dimensionalität der Datenmannigfaltigkeit.
  \item \textbf{PCA -- erklärte Varianz} -- die Anzahl der Hauptkomponenten, die 95\% der Gesamtvarianz erklären. Dieses Maß erfasst die globale lineare Dimensionalität.
\end{itemize}
%
Zwei Schätzer statt eines sind notwendig, weil sie verschiedene Aspekte messen:
PCA erfasst globale lineare Redundanz, TwoNN erfasst lokale Mannigfaltigkeitsstruktur.
Stimmen beide überein, ist die Aussage robust.
Divergieren sie, deutet dies auf nichtlineare Struktur hin, die PCA nicht erkennt.
%
\section{Phase~2: Distanz- und Fehleranalyse}
\label{sec:methodik_phase2}
%
Diese Phase beantwortet die Kernfrage:
Bleiben die geometrischen Beziehungen zwischen Dokumenten unter Quantisierung erhalten?
%
\subsection{Paarweise Distanzkorrelation}
%
Aus dem Dokumentenkorpus werden 10\,000 Paare gleichverteilt zufällig gezogen.
Für jedes Paar werden drei Distanzwerte berechnet:
Cosine-Similarity im Float-Raum, Hamming-Distanz im binären Raum und die entsprechende Distanz im TurboQuant-Raum.
%
Die Stichprobengröße von 10\,000 Paaren ist so gewählt, dass das 95\%-Konfidenzintervall für den Pearson-Korrelationskoeffizienten bei $r = 0{,}9$ eine Breite von weniger als $\pm 0{,}005$ besitzt.
Damit sind selbst kleine Unterschiede zwischen den Verfahren statistisch auflösbar.
%
Die Distanzerhaltung wird über zwei Korrelationsmaße quantifiziert:
%
\paragraph{Pearson~$r$.}
Der Pearson-Korrelationskoeffizient misst den linearen Zusammenhang zwischen Float-Distanz und komprimierter Distanz.
Ein hoher Wert nahe~1 bedeutet, dass die Quantisierung die Distanzen lediglich linear transformiert, ohne die geometrische Struktur zu verzerren.
%
\paragraph{Spearman~$\rho$.}
Der Spearman-Rangkorrelationskoeffizient misst, ob die Rangordnung der Distanzen erhalten bleibt.
Für Retrieval-Anwendungen ist dieses Maß entscheidender als Pearson~$r$:
Solange das Ranking korrekt ist, sind absolute Distanzverzerrungen irrelevant.
%
Die Kombination beider Maße erlaubt eine differenzierte Aussage:
Sind beide hoch, ist die Quantisierung strukturtreu.
Ist nur Spearman hoch, eignet sich die Repräsentation für Ranking, aber nicht für distanzbasierte Schwellwertentscheidungen.
Sind beide niedrig, ist die semantische Struktur substantiell beschädigt.
%
\subsection{Distortion}
%
Zusätzlich zur Korrelation wird die absolute Verzerrung gemessen.
Für jedes der 10\,000 Paare wird der Fehler $|d_\text{float} - d_\text{quantized}|$ berechnet, wobei die Distanzen zuvor auf $[0, 1]$ normiert werden.
%
Als Aggregatmaße dienen:
\begin{itemize}
  \item \textbf{MAE} (Mean Absolute Error) -- die durchschnittliche Abweichung, robust gegenüber Ausreißern.
  \item \textbf{RMSE} (Root Mean Squared Error) -- gewichtet große Abweichungen stärker und zeigt, ob einzelne Paare besonders stark verzerrt werden.
\end{itemize}
%
Während Pearson und Spearman die \emph{relative} Struktur bewerten, quantifizieren MAE und RMSE die \emph{absolute} Verzerrung.
Beide Perspektiven sind notwendig:
Ein Verfahren mit hoher Korrelation, aber hohem MAE transformiert die Distanzen zwar monoton korrekt, verschiebt aber das gesamte Distanzniveau -- was für Schwellwert-basierte Systeme problematisch wäre.
%
\section{Phase~3: Nachbarschaftserhaltung}
\label{sec:methodik_phase3}
%
Die Distanzkorrelation aus Phase~2 betrachtet globale Paare.
Für Retrieval ist jedoch die \emph{lokale} Struktur entscheidend:
Es kommt nicht darauf an, ob weit entfernte Dokumente korrekt als weit entfernt erkannt werden, sondern ob die nächsten Nachbarn eines Dokuments dieselben bleiben.
%
\subsection{Neighborhood Overlap}
%
Für jedes Dokument werden die $k = 100$ nächsten Nachbarn im Float-Raum bestimmt.
Anschließend werden die $k = 100$ nächsten Nachbarn im komprimierten Raum bestimmt.
Der Overlap ist der Anteil der Nachbarn, die in beiden Mengen enthalten sind:
\begin{equation}
  \text{Overlap}(i) = \frac{|N_k^\text{float}(i) \cap N_k^\text{quant}(i)|}{k}
\end{equation}
%
Dieser Wert wird über alle Dokumente gemittelt und als Funktion von $k$ (für $k \in \{10, 50, 100\}$) berichtet.
Die Variation von $k$ ist wichtig, weil Retrieval-Systeme typischerweise nur die Top-10 oder Top-100 verwenden -- der Overlap bei $k = 10$ ist für die Praxis relevanter als bei $k = 100$, aber statistisch instabiler.
%
\subsection{Trustworthiness}
%
Neighborhood Overlap ist symmetrisch und unterscheidet nicht zwischen zwei verschiedenen Fehlertypen.
Trustworthiness \cite{venna_2006_trustworthiness} differenziert:
Sie misst, ob im komprimierten Raum \emph{falsche} Nachbarn eingeführt werden -- also Dokumente, die im Float-Raum weit entfernt sind, im Bit-Raum aber nahe erscheinen.
%
Formal:
\begin{equation}
  T(k) = 1 - \frac{2}{Nk(2N - 3k - 1)} \sum_{i=1}^{N} \sum_{j \in U_k(i)} (r(i,j) - k)
\end{equation}
%
wobei $U_k(i)$ die Menge der Punkte bezeichnet, die im komprimierten Raum unter den $k$ Nachbarn von $i$ sind, aber im Float-Raum nicht, und $r(i,j)$ der Rang von $j$ bezüglich $i$ im Float-Raum ist.
%
Dieser Fehlertyp ist für Retrieval besonders schädlich:
Falsche Nachbarn bedeuten irrelevante Suchergebnisse, die dem Nutzer angezeigt werden.
Ein Wert von $T = 1$ bedeutet, dass keine falschen Nachbarn eingeführt werden.
%
Die Wahl, Trustworthiness statt des Gegenstücks Continuity zu verwenden, ist bewusst:
Continuity misst, ob echte Nachbarn verloren gehen.
Für die Forschungsfrage ist aber Trustworthiness aussagekräftiger, weil sie direkt die Precision der Nachbarschaft quantifiziert -- und Precision ist im Retrieval-Kontext die nutzerrelevante Größe.
Continuity wird nicht separat berichtet, da sie über den Neighborhood Overlap bereits implizit erfasst wird.
%
\section{Phase~4: Retrieval-Evaluation}
\label{sec:methodik_phase4}
%
Die bisherigen Phasen messen geometrische Strukturerhaltung.
Phase~4 prüft, ob sich diese geometrischen Beobachtungen in tatsächlicher Suchqualität niederschlagen.
%
\subsection{Retrieval-Protokoll}
%
Für jede Query im BEIR-Testset wird eine Suche in allen drei Repräsentationsräumen durchgeführt:
\begin{itemize}
  \item Float-Raum: Ranking nach Cosine-Similarity.
  \item Binärer Raum: Ranking nach Hamming-Distanz.
  \item TurboQuant-Raum: Ranking nach der entsprechenden Distanzfunktion.
\end{itemize}
%
Die Suche erfolgt exakt (Brute-Force), nicht approximativ.
Der Grund ist, dass approximative Verfahren (z.\,B. HNSW) eigene Fehler einführen, die den Quantisierungseffekt überlagern würden.
Das Experiment soll ausschließlich den Effekt der Kompression isolieren, nicht den einer approximativen Suche.
%
\subsection{Metriken}
%
Die Retrieval-Qualität wird über vier Standardmetriken bewertet:
\begin{itemize}
  \item \textbf{NDCG@10} -- die primäre Metrik, da sie sowohl Relevanzgrade als auch Rankingpositionen berücksichtigt und in der BEIR-Literatur als Hauptmetrik etabliert ist.
  \item \textbf{Recall@10} und \textbf{Recall@100} -- messen, welcher Anteil der relevanten Dokumente in den Top-$k$ enthalten ist. Recall@100 ist besonders relevant für zweistufige Systeme, in denen die erste Stufe eine Kandidatenmenge erzeugt.
  \item \textbf{MRR} (Mean Reciprocal Rank) -- misst, an welcher Position das erste relevante Dokument erscheint.
\end{itemize}
%
\subsection{Statistische Signifikanz}
%
Retrieval-Metriken variieren stark zwischen Queries.
Ein Unterschied in der mittleren NDCG@10 zwischen zwei Verfahren ist nur dann aussagekräftig, wenn er nicht durch zufällige Query-Schwankungen erklärbar ist.
%
Für jeden paarweisen Vergleich (Float vs. Binär, Float vs. TurboQuant, Binär vs. TurboQuant) wird daher ein Wilcoxon-Signed-Rank-Test auf den per-Query-Metriken durchgeführt.
Die Wahl fällt auf Wilcoxon statt eines gepaarten t-Tests, weil Retrieval-Metriken typischerweise nicht normalverteilt sind und Wilcoxon keine Verteilungsannahme erfordert.
%
Das Signifikanzniveau wird auf $\alpha = 0{,}05$ festgelegt.
Bei drei paarweisen Vergleichen pro Metrik wird eine Bonferroni-Korrektur angewandt ($\alpha_\text{korr} = 0{,}017$), um das multiple Testproblem zu berücksichtigen.
%
\section{Phase~5: Praktische Evaluation}
\label{sec:methodik_phase5}
%
Die bisherigen Phasen bewerten die Qualität der Kompression.
Phase~5 ergänzt die Perspektive der praktischen Anwendbarkeit.
%
\subsection{Speicheranalyse}
%
Für jede Repräsentation wird der theoretische Speicherbedarf pro Vektor berechnet:
\begin{itemize}
  \item Float32: $768 \times 32 = 24\,576$ Bit $= 3\,072$ Byte
  \item Binär: $768 \times 1 = 768$ Bit $= 96$ Byte
  \item TurboQuant: $768 \times 2 = 1\,536$ Bit $= 192$ Byte
\end{itemize}
%
Zusätzlich wird der tatsächliche Speicherverbrauch des gesamten Index gemessen, da Overhead durch Metadaten, Alignment oder Indexstrukturen den theoretischen Wert übersteigen kann.
%
\subsection{Laufzeitanalyse}
%
Die Suchgeschwindigkeit wird für verschiedene Korpusgrößen gemessen, um das Skalierungsverhalten zu erfassen.
Gemessen werden:
\begin{itemize}
  \item \textbf{Indexierungszeit} -- die Zeit, um alle Dokumente in die jeweilige Repräsentation zu überführen und zu indexieren.
  \item \textbf{Query-Latenz} -- die mittlere Zeit pro Query bei Brute-Force-Suche.
  \item \textbf{Throughput} -- Queries pro Sekunde.
\end{itemize}
%
Die Messung erfolgt auf den drei BEIR-Subsets, die unterschiedliche Korpusgrößen abdecken (SciFact: $\sim$5\,000, FiQA: $\sim$57\,000, TREC-COVID: $\sim$171\,000 Dokumente).
Damit ergibt sich ein natürliches Skalierungsspektrum, ohne synthetische Daten erzeugen zu müssen.
%
\section{Synthese und Interpretation}
%
Die abschließende Analyse verknüpft die Ergebnisse aller Phasen zu einer kohärenten Beantwortung der Forschungsfrage.
%
Konkret werden drei Leitfragen beantwortet:
\begin{enumerate}
  \item \textbf{Was bleibt erhalten?} Welche Aspekte der semantischen Struktur (globale Distanzen, lokale Nachbarschaften, Cluster) überleben die Quantisierung intakt?
  \item \textbf{Was geht verloren?} Wo treten systematische Verzerrungen auf, und lassen sich diese aus den in Phase~1 identifizierten Raumeigenschaften erklären?
  \item \textbf{Was bedeutet das für die Praxis?} Ab welchem Qualitätsverlust wird der Speicher- und Geschwindigkeitsvorteil durch unbrauchbare Suchergebnisse zunichte gemacht?
\end{enumerate}
%
Diese Struktur stellt sicher, dass die Arbeit nicht bei der Beobachtung stehen bleibt, sondern eine wissenschaftlich fundierte Empfehlung abgeben kann, unter welchen Bedingungen welche Kompressionstiefe gerechtfertigt ist.
